CRDT 是 conflict-free replicated data type 的缩写，通常译为无冲突复制数据类型。它为一个复制对象规定更新及并发合并语义，使收到同一组更新的副本按确定规则得到同一状态~\cite{preguica2018crdt}。先从一个只增长计数器看这些规则怎样工作。

\subsection*{用来源累计值识别重复传播}

先为已确认事件做一项计数，计数本身不触发付款。两个来源 A、B 的初始分量都为 0，状态写成向量 $p=(p_A,p_B)$，显示值为两个分量之和。A 每确认一项新的计数事件，只增加自己的分量，B 同样只写自己的分量。这样的只增长计数器通常称为 G-counter，G 表示 grow，即增长。

A 生成第一个增量以后，状态为 $(1,0)$。它把这份状态发送给 B，B 对每个来源取自己与来件的较大累计值，得到 $(1,0)$。同一状态再到一次，每个分量的最大值都保持。B 随后生成自己的一个增量，变为 $(1,1)$，再与 A 同步，双方都显示 2。

为什么不能把两个副本的显示值相加？B 的 1 已经包含 A 的第一项增量，A 的 1 也来自同一项。直接相加会把它数两次。来源分量说明哪些累计值覆盖同一串更新，每分量取最大才保留较长的同源前缀。

\CThirteenFigure{ch13-original05}{G-counter 对同来源累计值取最大，同一份状态重传保持一次增量}{fig:ch13-gcounter}

来源身份与累计数因此进入持久状态。A 重启后若把自己的分量重置为 0，再为一个新事件生成 $(1,0)$，已有副本会把它与旧 $(1,0)$ 合并为 1，新事件就被遗漏。另一个错误是两个活跃进程都以 A 的身份从 0 产生第一个增量，二者合并也只有 1。每个来源独占并串行分配自己的计数，重启以后从已经保存的累计值继续。

重启时分配新来源身份可以避免重用旧计数，但状态中会多出一个分量。设备不断加入、离线和更换身份时，来源数量随之增长。合并省去了逐次同步报文记录，状态中增加了来源登记与每源累计值。旧来源退出也要有处理规则：删除 A 的分量以后，若保存旧 A 状态的节点还能直接回来，接收端就要再次解释它带回的旧信息。

\subsection*{用半格说明为什么合并能够收敛}

刚才的两个分量解决了一个具体问题：A 已经数过的第一项事件，到 B 以后又经第三个副本转回来，怎样只保留一次？逐来源取最大值把同源累计前缀合在一起，同时保留不同来源的增量。接下来把这份关系写清楚，就能判断它在更多副本、多轮同步中为什么成立。

对计数向量定义 $p\preceq q$，当且仅当每个来源 i 都满足 $p_i\le q_i$。若 $p=(1,0)$、$q=(0,1)$，双方互不包含。它们最小的共同上界是 $(1,1)$：它至少包含双方每个分量，又没有额外增加分量。

任意两个状态都有唯一最小共同上界的偏序结构，称为 join 半格。这里的 join 指最小上界合并，记为 $p\sqcup q$。符号 $\sqcup$ 在本章总表示这项合并运算。对 G-counter，它就是逐分量取最大，对只添加集合，它就是集合并。

状态型 CRDT 将状态发给其他副本，接收端以 join 合并来件。它的基本构造要求状态按某个信息偏序形成 join 半格，本地更新沿偏序膨胀，合并取两个状态的最小共同上界~\cite{preguica2018crdt}。\footnote{Preguiça、Baquero 与 Shapiro 的作者综述分别讨论并发语义与同步模型。本节采用其中状态型构造的条件，来源累计值、购买入口与故障恢复的具体历史是本书的展开。}
膨胀保证本地演进保留已有信息，join 则把不同副本的已知信息合到共同上方。

为什么 join 天然具有前面的三项性质？两个状态的最小共同上界与输入次序无关，所以交换。一个状态与自己合并，最小上界就是自己，所以幂等。三个状态无论先合并哪两个，最后都是包含三者的最小上界，所以结合。重复、重排和分组由同一结构分别处理。

从初始最小状态开始，每次合法更新都留下更大状态，每次合并又保留双方信息，已有更新因而可以沿同步路径继续传播。对一组停止新增的更新，只要承载它们的状态有可恢复的副本，网络最终允许反复同步，并且每份更新最终到达所有目标副本，各处便汇到这一组更新的同一 join。更新持续产生时，对任何固定有限前缀也可以按相同传播条件讨论最终纳入~\cite{preguica2018crdt}。

\CThirteenFigure{ch13-original06}{本地膨胀保存已有信息，join 积累不同分支，反复同步让更新逐步到达各副本}{fig:ch13-join}

这份推导同时给出崩溃恢复的要求。若一个尚未传播的新分量只存在于丢失的内存，其他副本没有材料把它合回来。若恢复把来源计数退回过去，又让同一来源继续分配，标记唯一性和累计覆盖关系都会改变。对象持久化保存已经形成的更新，来源分配延续原计数，重新加入流程再把恢复所得状态接回同步路径。

允许减少显示值时，可以保存两个只增长向量 p、n，分别累计增加与减少。合并时各自逐分量取最大，显示值为
\[
  \operatorname{value}(p,n)=\sum_i p_i-\sum_i n_i.
\]
这种正负计数器称为 PN-counter。增加 $n_i$ 以后，信息状态膨胀，显示值下降。例如一项收款增加 $p_A$ 3000，随后一项退款增加 $n_A$ 1000，累计分量分别为 3000、1000，显示净收款 2000。若减少量累计到了 4000，结构就表达 $-1000$，允许生成多少减少量由余额下限的业务规则决定。

\CThirteenFigure{ch13-pn-observation}{两个累计分量继续增长，净收款从 3000 变为 2000}{fig:ch13-pn-observation}

\subsection*{入口新建增量与重复发送旧状态怎样分开}

现在把计数器接到付款入口。共同初态为支付平台余额 10000，G-counter 为 $(0,0)$，S1 尚未受理。错误入口每收到一次 S1，就为本次请求新增一项付款计数，另建一份支付记录，并用新的平台键安排一次 3000 分支付。

第一次来件让 A 从 $(0,0)$ 变成 $(1,0)$。入口创建 C1/\code{pay-C1}，支付形成 X1，响应丢失。S1 重试以后，入口又生成第二个增量，A 变成 $(2,0)$，并创建 C2/\code{pay-C2}，平台可以形成 X2。此后把 $(2,0)$ 重传一百次，所有副本都会稳定得到 2，合并准确保留了入口产生的两项增量。

入口若先按 S1 取得原 C1，已经受理就恢复原关联，只有首次受理分支产生一项增量，后续状态传播便可以用 G-counter 压缩。还有一种故障只影响统计：入口多加了一次计数，实际支付却继续沿同一 \code{pay-C1}，且每次可能生效的到达都受平台身份保护。在这条历史中，统计为 2，平台可以只有一项 X1，余额为 7000。

\CThirteenFigure{ch13-original07}{入口若因重试另建 C2 和平台键，后续合并会稳定保留两项增量与付款}{fig:ch13-ingress}

把计数事件改为稳定业务事实，也有另一条路。例如只统计不同 C 的已核验成功，以 C1/X1 的关系作为成员依据，再求集合大小。重复加入同一 C1 不增加计数，两项不同的成功 C1、C2 则都留下成员。空间成本转为成功事实身份，支付入口则按 S 与固定平台关系形成预期的 C/X 关联。

\subsection*{操作型传播需要说明交付什么}

状态型发送累计状态，操作型 CRDT 则传播已经形成的操作效果。源副本先根据本地状态准备一份效应器，再在源和其他副本应用它。效应器是编码具体状态变化的数据或可执行描述，例如增加 1，或者删除源端已经观察到的这些标记。准备过程本身只计算这份描述，持久登记与应用由系统安排~\cite{preguica2018crdt}。

以加法计数器为例，两项不同效应器 \code{inc(1)}、\code{inc(1)} 交换，按任意顺序各应用一次都得到 2。若其中一份被重复应用，结果就为 3。常见操作型构造把每份效应器在各副本应用一次作为传输层提供的前提，通信实现可以通过可靠投递加消息身份去重实现它。底层允许重复到达时，接收端按原消息身份识别已经应用的效果，或直接采用自身幂等的效应器。

传输顺序也影响设计。在因果交付中，若更新 $u_2$ 在源端已经观察到 $u_1$，接收端就先交付 $u_1$，再交付 $u_2$。只要求并发效应器交换的构造依赖这种顺序。若完全不保留因果顺序，允许重排的全部效应器都要满足所需交换条件~\cite{preguica2018crdt}。

例如源端看见标签 $(\mathrm{gift},A1)$ 后删除 gift，准备的效应器明确写移除 A1。传输若保证 A1 的添加先到、删除后到，并且各应用一次，接收端可以直接移除。若删除先到，接收端删除空项并丢弃这份信息，迟到的 A1 又会出现。保存 A1 的删除标记可以处理这条乱序历史，接收端因此多保存一份删除依据。

再把两类报文放到同一条故障历史中。A 已经持久生成第 7 项计数，接收方 B 已经处理，响应丢失。累计状态的重发内容为 \code{(A,total=7)}，B 原来保存的 A 分量也是 7，取最大后保持 7。若重发内容为第 7 份 \code{inc(1)}，裸加法会从 7 变成 8。接收端按效应器身份 \code{(A,7)} 找到已应用依据，才能跳过这次重复。

\CThirteenFigure{ch13-state-effector}{响应丢失后，重发累计值与重发加一效果，接收端采用不同的比较规则}{fig:ch13-state-effector}

A 重启以后，从持久计数 7 继续生成下一项。累计状态路径发送新的 \code{(A,total=8)}，B 将分量推进到 8。操作路径生成新的效应器身份 \code{(A,8)}，B 将它与原 \code{(A,7)} 区分，并应用这一项新效果。重发旧内容与产生新内容在两条路径中都有固定依据，进程重启也沿这些依据继续。
\FloatBarrier
